\documentclass[11pt]{amsart}
\usepackage[T1]{fontenc}
\usepackage{lmodern,amsmath,amssymb,amsthm,mathtools,microtype,booktabs,array}
\usepackage[margin=1.12in]{geometry}
\usepackage[hidelinks]{hyperref}
\newtheorem{theorem}{Theorem}[section]
\newtheorem{proposition}[theorem]{Proposition}
\newtheorem{lemma}[theorem]{Lemma}
\newtheorem{corollary}[theorem]{Corollary}
\theoremstyle{definition}\newtheorem{definition}[theorem]{Definition}
\theoremstyle{remark}\newtheorem{remark}[theorem]{Remark}
\newcommand{\TV}{\operatorname{TV}}
\newcommand{\Ber}{\operatorname{Ber}}
\newcommand{\E}{\mathbb E}
\newcommand{\Pp}{\mathbb P}
\newcommand{\1}{\mathbf 1}
\newcommand{\cP}{\mathcal P}
\newcommand{\cA}{\mathcal A}
\newcommand{\cR}{\mathcal R}
\newcommand{\osc}{\operatorname{osc}}
\newcommand{\supp}{\operatorname{supp}}
\newcommand{\conv}{\operatorname{conv}}
\newcommand{\Qstar}{Q_*}
\title[Controlled statistical minimax]{General Theta Foundations I:\\
Controlled Experiment Duality and Exact Marked Minimax Laws}
\author{Qian Qi}
\date{September 24, 2026. Revision 25}
\subjclass[2020]{Primary 62B05, 62C20; Secondary 60J10, 68Q45, 93E20}
\keywords{Statistical experiments, causal realization, predictive distributions,
finite memory, minimax duality, physical certification}
\emergencystretch=2em
\begin{document}
\begin{abstract}
We study finite-resource statistical comparison when training operations
are chosen causally and a validation test is frozen before fresh candidate
and target preparations. For finite controlled experiments with a compact
parameter space, we prove an architecture-uniform posterior-predictive
dual in which action choice remains inside a Bellman recursion and the
adversarial prior is chosen only once. The actions need not be garblings
of a common experiment. We solve the one-preparation minimax problem
for a continuous family of marked product experiments, identifying both
an optimal response and a two-point least-favorable prior by an explicit
sum-of-squares identity. For the collision target the exact value is
$(85-7\sqrt{73})/64$. Combined with a deterministic two-preparation audit,
this gives an enlarged exact preparation segment at every peak width
at least twelve. The microscopic collision and positive-noise bridge
are proved within the article. A separate revelation family has the
exact joint law $(1-(1-\alpha)^N)(1-1/K)$ for preparations and decision
memory. Compatible finite-architecture certificates and a priced
classical autonomous comparison complete the resource analysis.
\end{abstract}
\maketitle
\tableofcontents

\section{Introduction}
A statistical experiment determines not only which decisions can be
made, but which observations and retained continuations are required to
make them. In a controlled experiment, an action may change the
information available at the next preparation. In a finite-memory
experiment, retaining one continuation may preclude retaining another.
A resource theorem must distinguish these two constraints rather than
replace both by optimization within one chosen execution graph.

We consider an unknown parameter fixed throughout an audit. Training
preparations are conditionally independent draws from causally selected
experimental rows. The auditor then freezes a validation design and an
event and compares fresh candidate and target draws through that event.
All persistent randomness, report buffers and validation accumulators
belong to the register. This model separates the simulator's resources
from those of its auditor and keeps the mark of the actual preparation
inside the statistical comparison.

The first result is a controlled posterior-predictive dual
(Theorem~\ref{thm:controlled}). Its state is an unnormalized posterior
measure together with the legal protocol state. At the terminal step,
optimization over an event gives a positive-part functional, equivalently
a posterior-predictive total variation. Earlier steps optimize the
causal experimental action. One minimization over initial priors yields
the exact unrestricted-memory minimax value. Every such prior, or a
Bellman supersolution, gives a converse uniform over admissible
finite-memory architectures. The common-dominating-row theorem becomes
a specialization in which the experimental choice can be removed.
The controlled formulation also applies when the available actions are
statistically incomparable. A strict adaptive advantage persists under
full-support perturbations (Proposition~\ref{prop:adaptive-gap}).

Our principal explicit minimax calculation concerns the product family
$b(p,q)=\Ber(p)\otimes\Ber(q)\otimes\Ber(1/2)$ and targets
$Q_\gamma$ on the four marked diagonal atoms. For
$0\le\gamma\le1-1/\sqrt2$, put
\[
 u_\gamma=\gamma-2+\sqrt{\gamma^2-2\gamma+5}.
\]
Theorem~\ref{thm:marked-exact} gives
\[
                 U_1(b,Q_\gamma)=(3+u_\gamma^2)/8.
\]
The optimal response and least-favorable prior are explicit. Their
matching is certified over the entire parameter square by a square
identity, not by an optimization grid. The endpoint of the displayed
$\gamma$ interval is the point where a posterior response coefficient
changes sign. The zero-preparation value is also exact. At the physical
mark bias $\gamma=1/4$ the one-preparation value is
\[
                   u_{\rm ex}=(85-7\sqrt{73})/64.
\]
This resolves the optimization problem hidden behind the earlier
rational one-preparation certificate.

The physical consequence uses the same marked collision experiment,
not a surrogate alphabet with a different cost convention. Define
\begin{equation}\label{eq:constants}
 u_1=\frac{4035777}{10240000},\qquad
 m_2=\frac{12-3\sqrt3}{16},\qquad \beta_0=\frac1{300}.
\end{equation}
The older rational bound $u_1$ remains a valid explicit certificate.
The exact value $u_{\rm ex}$ improves it. We prove
\begin{equation}\label{eq:headline}
 N_c^{\rm phys}(W,Q)=2\quad\text{for}\quad
 W\ge12,\quad\TV(Q,Q_*)\le\beta,\quad
 u_{\rm ex}+\beta\le c<m_2-\beta.
\end{equation}
Here $N$ counts candidate training preparations, and one candidate and
one target validation are charged in addition. The upper construction
is deterministic and has the complete bit-level profile
$(1,2,3,3,4,5,5,10,12,10,10,7,3)$.
Definition~\ref{def:physical-class} and Theorem~\ref{thm:local-physical}
prove locally every interface and physical approximation fact used in
this deduction. Equation~\eqref{eq:headline} is sharp in preparation
count on its stated width range; it does not assert width optimality
at two preparations or an exact value of $U_2$.

A different theorem determines both coordinates in a revelation family.
With reveal probability $\alpha$, $N$ preparations and $K$ retained
decision states, Theorem~\ref{thm:joint-law} proves
\[
       V_N^{\rm er}(K)=(1-(1-\alpha)^N)(1-1/K).
\]
The converse allows every stochastic encoder and legal stopping rule;
the matching machine requires no uncharged revelation flag. The formula
yields exact strict thresholds and matching sample and state scales.
The memory coordinate here is a decision-cut width, not the collision
machine's bit-level peak. Keeping those coordinates distinct is part
of the resource theorem.

For a specified finite architecture, compatible occupation arrays retain
all common-row constraints, including autonomous row reuse. The
subdivision theorem gives an explicit $HS/m$ certification error with
controlled row-block degree. These are effective finite-architecture
certificates, not a sharp universal complexity law. We also retain a
resource-accounted autonomous binary comparison whose stationary formula
is classical. None of these auxiliary results replaces the exact
minimax or physical arguments.

\subsection{Relation to earlier work}
The comparison and minimax mechanisms are classical
\cite{Blackwell,Sion}. Finite-horizon belief-state control and
piecewise-linear value functions are established in
\cite{SmallwoodSondik}. Our controlled dual uses these mechanisms with a
fixed adversarial parameter, frozen validation and explicit causal
resource accounting. It is not a claim to have invented belief-state
control or the minimax theorem. The exact marked-family saddle points,
their equality certificate and their physical use are stated separately
from these general mechanisms.

Polynomial stochastic-policy formulations and positivity certificates
have substantial antecedents \cite{MM,Putinar}. Finite-memory testing,
randomization and time-invariant rules are treated in
\cite{HC,HCrandom,FH,CFH}. The stationary formula in
Section~\ref{sec:autonomous} is explicitly attributed to Hellman and
Cover. The May 2026 work of Managoli and Prabhakaran \cite{MP} treats
memory-constrained adversarial binary hypothesis testing, including
randomized time-invariant machines and adaptive distribution choices.
Its objective and adversarial mechanism differ from the fixed-parameter,
training--validation score here. No general priority for randomized
state savings is asserted.

Norberg's filtered comparison \cite{Norberg} belongs to the surrounding
comparison theory. We make no non-overlap or originality claim about
that general theory without a proof-level comparison of the original
paper. The novelty claims under consideration are the explicit theorems
and resource realizations proved here, not an asserted absence of
filtered-experiment predecessors.

\subsection{Organization and preserved development}
We first fix the architecture and common-row framework, then prove the
controlled dual and exact marked minimax law. The physical bridge and
charged two-preparation construction give \eqref{eq:headline}. The
joint revelation law, finite-architecture certification and autonomous
comparison follow. Complete mathematical volumes retain the unaltered
predecessor article and development, including the general continuation,
Gaussian transport and historical pipeline proofs \cite{Companion}.
These remain mathematical content, not substitutes for a proof in the
present article. The independent A2 geometric chain and the original
B4 and broad C2 obligations are not claimed to follow from the finite
statistical theorems alone.

\section{Experiments, architectures and attainable resources}
\label{sec:model}
Throughout, total variation is
$\TV(P,Q)=\sup_{0\le h\le1}(P-Q)h=\frac12\sum_x|P(x)-Q(x)|$.
An experimental parameter is unknown and fixed throughout an execution.
Controller randomization is independent of that parameter. A reset
preparation is a fresh draw, independent of earlier preparations given
the parameter. A controller may not recover discarded information from
an uncharged tape, clock, selector or calibration table.

\begin{definition}[Finite architecture]
A finite architecture $A$ specifies a finite causal execution graph, its
legal experimental operations, the observable inputs to each controller
row, and a register alphabet at every cut. It also specifies which row
occurrences share a stochastic table. A clocked architecture may use
different rows at different scheduled times. An autonomous architecture
may use time only through a charged state. For a parameter family $\Theta$
and terminal reward in $[-1,1]$, let
\[
 V(A)=\sup_{q\in\prod_r\Delta_{s_r}}\inf_{\theta\in\Theta}
          g_\theta(A,q).
\]
If $\cA(R)$ is the class of architectures obeying a global resource budget
$R$, its outer value is
\[
 V^{\rm out}(R)=\sup_{A\in\cA(R)}V(A).
\]
Both the architecture and the simulator parameter are fixed before the
controller's random choices. In particular, a mixture of complete
controllers is not automatically an admissible controller at the original
register profile.
\end{definition}

For a finite hidden execution graph, compatibility can be stated without
division. Let $x_{\theta i r}$ be the incoming mass of hidden history $i$
at observable row $r$, and let $f_{\theta i r j}$ be the outgoing mass
through choice $j$. Include the usual fixed-kernel flow equations, as
well as
\begin{equation}\label{eq:occupation}
 f_{\theta i r j}\ge0,\qquad
 \sum_jf_{\theta i r j}=x_{\theta i r},\qquad
 f_{\theta i r j}x_{\theta'i'r}
 =f_{\theta'i'rj}x_{\theta i r}.
\end{equation}
For shared rows, the indices $i,i'$ include all their occurrences.

\begin{proposition}[Compatible realization]\label{prop:realization}
For a specified finite architecture, nonnegative flow arrays satisfying
\eqref{eq:occupation} are exactly those realized by a common stochastic
row at every declared observable input, including at zero incoming mass.
\end{proposition}
\begin{proof}
Necessity follows from $f_{\theta i rj}=x_{\theta i r}q_r(j)$.
For sufficiency, if some incoming mass at row $r$ is positive, define
$q_r(j)$ by dividing the corresponding outgoing mass by that incoming
mass. Conservation makes $q_r$ stochastic, and the cross-products force
the same ratio at every other positive incoming mass. If an incoming
mass vanishes, nonnegativity and conservation force all its outgoing
masses to vanish. If the entire row is unreachable, choose any legal
stochastic row. Induction through the finite execution graph reconstructs
the prescribed flows. The same ratio argument applies to identified
occurrences of an autonomous row.
\end{proof}

This formulation is inherited from the complete companion and is included
to fix the controller class. Compact semialgebraic feasibility and an
Archimedean positivity theorem provide a complete semialgebraic certificate
hierarchy for a \emph{fixed finite architecture}. They do not interchange
the two optimizations in the definition. Section~\ref{sec:quantitative}
gives a quantitative alternative for the inner problem. The next section
supplies a different certificate that survives the outer optimization.

\section{A statistical dual that is uniform over architectures}
\label{sec:dual}
Let $\Theta$ be compact, let $\Omega$ be finite, and let
$\theta\mapsto P_\theta\in\cP(\Omega)$ be continuous. Fix
$Q\in\cP(\Omega)$. An $N$-preparation reset selector observes
$Z=(Z_1,\ldots,Z_N)\sim P_\theta^{\otimes N}$ and chooses a response
$h_Z\in[0,1]^\Omega$. Its expected validation score is
\begin{equation}\label{eq:selector}
 G_\theta(h)=\sum_{z\in\Omega^N}P_\theta^{\otimes N}(z)
                           (Q-P_\theta)h_z.
\end{equation}
The response is fixed before two independent validation draws. Operationally
it can be realized by choosing a random event with those coordinate
marginals and retaining that event for both validations. No memory bound
is imposed in the definition of
$U_N(P,Q)=\max_h\min_\theta G_\theta(h)$.

\begin{theorem}[Posterior-predictive dual]\label{thm:dual}
Both extrema in the following identity are attained:
\begin{align}
 U_N(P,Q)
 &=\min_{\mu\in\cP(\Theta)}
   \sum_{z,\omega}
   \left[\int P_\theta^{\otimes N}(z)
               \{Q(\omega)-P_\theta(\omega)\}\,d\mu(\theta)\right]_+
                                                        \label{eq:dual}\\
 &=\min_{\mu\in\cP(\Theta)}
   \sum_z m_\mu(z)\TV(Q,\overline P_\mu(\cdot\mid z)).
                                                        \label{eq:predictive}
\end{align}
Here $m_\mu(z)=\int P_\theta^{\otimes N}(z)d\mu(\theta)$ and, when
$m_\mu(z)>0$,
\[
 \overline P_\mu(\omega\mid z)
 =\frac{\int P_\theta^{\otimes N}(z)P_\theta(\omega)d\mu(\theta)}
        {m_\mu(z)}.
\]
A zero-mass term in \eqref{eq:predictive} is defined to be zero.
\end{theorem}
\begin{proof}
The response space is a finite-dimensional compact convex box.
Probability measures on the compact space $\Theta$ form a weakly compact
convex set. The integral of \eqref{eq:selector} is continuous and affine
in each argument. The minimum over parameters equals the minimum over
probability measures. The minimax theorem \cite{Sion} therefore interchanges
the maximum and minimum. For fixed $\mu$, each response coordinate is
independent in the maximization and its contribution is the positive part
of its coefficient, proving \eqref{eq:dual}. The coefficient vector at a
fixed $z$ has zero sum. For positive $m_\mu(z)$ it equals
$m_\mu(z)(Q-\overline P_\mu(\cdot\mid z))$, proving
\eqref{eq:predictive}. If $m_\mu(z)=0$, every nonnegative likelihood
integrand vanishes $\mu$-almost surely and all coefficients are zero.
Continuity on the two compact strategy spaces gives attainment.
\end{proof}

\begin{theorem}[Dominating-row converse]\label{thm:outer}
Suppose every permitted training row is a parameter-independent causal
garbling of one fresh draw from $P_\theta$. Suppose every permitted
validation row is the same garbling of the target and candidate laws,
and every selected row--event pair is fixed before validation. Then
$U_N(P,Q)$ bounds the worst-case expected score of every auditor using
at most $N$ candidate training preparations, regardless of its register
profile, feedback schedule, independent randomization, or stopping rule.
Independent target-only calibration, whose law does not depend on
$\theta$, does not invalidate the bound. If the all-on row and all its
events are permitted, the bound is attained in the class with no memory
cap and $N$ training preparations.
\end{theorem}
\begin{proof}
Generate $N$ independent all-on observations in advance. At the $t$th
requested preparation, apply the declared garbling to observation $t$,
using only previously simulated reports and independent randomization.
This simulates the adaptive transcript with its exact law. A stopped
execution ignores the unused observations; it does not reveal them to
the auditor. Pull back its final row--event response to the all-on
alphabet. Conditional on the entire generated training word and the
independent random seed, this is a function in $[0,1]^\Omega$.
Averaging the seed, including any target-only calibration, produces one
response $h_z$ for each full word $z$. The score is consequently
\eqref{eq:selector}. This proves the upper bound for the union of all
architectures. Conversely, retaining the entire word and then a selected
event implements any such finite response kernel when memory is
unrestricted. Fresh validation gives exactly \eqref{eq:selector}.
\end{proof}

\begin{remark}[What the dual does not convexify]
The measure $\mu$ is an adversarial prior used to bound the minimum over
\emph{fixed} parameters. Its likelihood is
$\int P_\theta^{\otimes N}d\mu(\theta)$, not
$(\int P_\theta d\mu(\theta))^{\otimes N}$. It neither enlarges the
private candidate class to mixtures nor gives the auditor a free
persistent selector. Attainment of $U_N$ without a memory cap does not
assert attainment at any specified finite profile.
\end{remark}

The state in \eqref{eq:predictive} is generated by the actual preparation
law, not by a geometric measure chosen independently of the experiment.
For sequential continuation one retains the posterior $\mu_z$ or its
future-test equivalence class; a single predictive marginal need not be
closed under further observations. This is the distinction between a
one-step response quotient and a future-predictive quotient.

\begin{corollary}[Finite moment support for the product family]
\label{cor:moments}
For two independent Bernoulli report bits with parameters $(p,q)$ and an
independent fair mark, an optimal prior in \eqref{eq:dual} can be chosen
with at most $(N+2)^2$ atoms. The two report counts $(k,\ell)$ are a
causally updatable sufficient statistic for computing the dual.
\end{corollary}
\begin{proof}
Conditional on the two counts, all arrangements of report bits and all
mark words have parameter-independent probabilities. Summing their
coefficients reduces \eqref{eq:dual} to the $(N+1)^2$ count pairs.
Their likelihoods are
\[
 {N\choose k}p^k(1-p)^{N-k}
 {N\choose\ell}q^\ell(1-q)^{N-\ell}.
\]
After multiplication by $Q(\omega)-P_{p,q}(\omega)$, every integrand
has degree at most $N+1$ in each variable. Thus the objective depends
only on the $(N+2)^2$ rectangular moments. Their constant coordinate is
one. Carath\'eodory's theorem in the remaining affine space replaces an
optimal moment vector by a convex combination of at most $(N+2)^2$
point evaluations. Counts update by addition of the current bits.
No claim that this count register is optimal for the decision is made.
\end{proof}

\section{Controlled experiments without a dominating action}
\label{sec:controlled}
The reduction to one dominating row is not needed for a statistical
dual, provided causal action choice is retained in the dual recursion.
We give the finite-action theorem with a compact parameter space and an
explicit protocol state. Its equality is an unrestricted-memory
statement; its upper certificates apply to every smaller-memory class.

Let $\Theta$ be compact metric and $E$ a finite set of protocol states.
At $e\in E$ let $\mathcal A_e$ be a finite set of training actions.
Action $a$ produces $x$ in a finite alphabet $X_a$ with law
$P^a_\theta$ and changes the protocol state to $T(e,a,x)$.
The functions $\theta\mapsto P^a_\theta(x)$ are continuous.
An unknown $\theta$ is fixed throughout the experiment. Conditional on
it, each requested preparation uses a fresh independent draw from the
selected row. Legal validation designs form a nonempty finite set
$\mathcal B_e$. Design $b$ has a finite outcome alphabet $Y_b$, target
law $Q^b$ and continuous candidate law $V^b_\theta$.
A policy may stop after at most $N$ training actions. It then freezes
$b\in\mathcal B_e$ and an event on $Y_b$, before one candidate and one
target validation. Randomization is independent of $\theta$.
No training action is assumed to dominate another.

For a finite nonnegative measure $\nu$ on $\Theta$ write
$\nu_{a,x}(d\theta)=P^a_\theta(x)\nu(d\theta)$ and set
\begin{align}
 \Phi_e(\nu)&=\max_{b\in\mathcal B_e}\sum_{y\in Y_b}
       \left[Q^b(y)\nu(\Theta)-\int V^b_\theta(y)d\nu\right]_+,
                                                   \label{eq:controlled-terminal}\\
 F_0(e,\nu)&=\Phi_e(\nu),\notag\\
 F_n(e,\nu)&=\max\left\{\Phi_e(\nu),\
        \max_{a\in\mathcal A_e}\sum_{x\in X_a}
           F_{n-1}(T(e,a,x),\nu_{a,x})\right\}.
                                                   \label{eq:controlled-recursion}
\end{align}
If no training action is legal, only the stopping term is retained.
All the functions vanish at the zero measure, so no posterior division
is required on null histories. For positive mass, \eqref{eq:controlled-terminal}
is the mass times the largest posterior-predictive total variation.

\begin{theorem}[Controlled posterior-predictive dual]
\label{thm:controlled}
Let $\mathcal U_N(e_0)$ be the maximum worst-case expected validation
score over all randomized causal policies in the declared protocol,
with at most $N$ training preparations and no memory cap. Then
\begin{equation}\label{eq:controlled-dual}
       \mathcal U_N(e_0)=\min_{\mu\in\cP(\Theta)}F_N(e_0,\mu).
\end{equation}
Both extrema are attained. Each $F_n(e,\cdot)$ is the maximum of
finitely many continuous linear functionals on the cone of measures;
in particular it is convex, continuous and positively homogeneous.
For every prior $\mu$, $F_N(e_0,\mu)$ bounds the worst-case score of
every finite-memory architecture implementing this protocol.
\end{theorem}
\begin{proof}
There are finitely many histories of length at most $N$ and finitely
many event masks at each stopping history. Consequently there are
finitely many deterministic policy trees, denoted by $\Pi_N(e_0)$.
For $\pi\in\Pi_N(e_0)$ its payoff $g_\pi(\theta)$ is a finite sum
of finite products of continuous primitive probabilities and is
continuous in $\theta$.

A randomized causal policy induces the same outcome laws as a probability
distribution on these trees: sample its finite random decision table in
advance and condition on it. Conversely, with no memory cap, a sampled
tree can be retained and followed. Thus the unrestricted payoff vectors
are exactly the convex hull of the $g_\pi$. This step is not an assertion
that mixing machines preserves a specified register width.

The simplex on $\Pi_N(e_0)$ and $\cP(\Theta)$ are compact and convex.
The expected payoff is continuous and bilinear. Minimax therefore gives
\[
 \max_{\lambda\in\Delta(\Pi_N)}\min_\theta
       \sum_\pi\lambda_\pi g_\pi(\theta)
 =\min_\mu\max_{\pi\in\Pi_N}\int g_\pi\,d\mu.
\]
For a fixed measure, the terminal maximum is
\eqref{eq:controlled-terminal}, because an event includes precisely the
positive coordinates of its signed predictive measure. At a training
node one selects an action and, after its observed outcome, selects a
continuation tree separately on that branch. The unnormalized branch
measure is exactly $\nu_{a,x}$. Backward induction proves
\eqref{eq:controlled-recursion} and identifies the right side with
\eqref{eq:controlled-dual}. It also proves the finite-linear-functional
representation. Compactness gives an optimizing prior and an optimizing
mixed tree. Finally, the minimum over fixed parameters of any policy is
at most its average under $\mu$, which is at most $F_N(e_0,\mu)$.
This remains true after restricting its memory or architecture.
\end{proof}

The minimization in \eqref{eq:controlled-dual} occurs once, outside the
recursion. It cannot be replaced by fresh adversarial choices after
observations. The posterior carries dependence on one fixed unknown
parameter. Also, a Bayes-optimal tree for a minimizing prior need not by
itself be a minimax policy; the saddle-point mixture is supplied by the
finite game in the proof.

\begin{corollary}[Bellman upper certificates]\label{cor:bellman-certificate}
Suppose functions $\Psi_n(e,\nu)$ satisfy $\Psi_0\ge\Phi$ and, for
$n\ge1$, satisfy
\[
 \Psi_n(e,\nu)\ge\Phi_e(\nu),\qquad
 \Psi_n(e,\nu)\ge\sum_x\Psi_{n-1}(T(e,a,x),\nu_{a,x})
                  \quad(a\in\mathcal A_e).
\]
Then $\Psi_N(e_0,\mu)$ is an upper certificate for every admissible
architecture and every prior $\mu$. The family $F_n$ is the least such
family. An optimal prior can be chosen with at most $d+1$ support points,
where $d$ is the affine dimension of
$\{(g_\pi(\theta))_{\pi\in\Pi_N}:\theta\in\Theta\}$.
\end{corollary}
\begin{proof}
Induction in \eqref{eq:controlled-recursion} gives $F_n\le\Psi_n$.
The equality case is the recursion itself. The objective depends on a
prior only through its vector of integrals of the finitely many
$g_\pi$. An optimal such vector belongs to the convex hull of their
compact evaluation image. Carath\'eodory's theorem gives the stated
support bound while preserving every pure-tree payoff.
\end{proof}

\begin{proposition}[Causal stability]\label{prop:controlled-stability}
Consider two models with the same protocol, alphabets and policy class.
Suppose their training rows differ uniformly in total variation by at
most $\epsilon$, their candidate validation rows by at most $\eta$,
and their target validation rows by at most $\beta$. Their unrestricted
values, or their values in any common specified memory class, differ by
at most $2N\epsilon+\eta+\beta$.
\end{proposition}
\begin{proof}
Use the same policy randomization and couple corresponding training
rows while the histories agree. The probability of a first discrepancy
is at most $N\epsilon$, even with feedback and stopping. Conditional
expected validation rewards lie in $[-1,1]$, so differing training
histories contribute at most $2N\epsilon$. On agreeing histories the
same frozen design and event are used, and their conditional mean
rewards differ by at most $\eta+\beta$. The bound is uniform in the
fixed parameter and in the policy. Taking minima and suprema preserves
it, with the reverse inequality obtained by interchanging the models.
\end{proof}

\subsection{An informative action choice}
The following calculation distinguishes the controlled theorem from a
common-row reduction. Let $\Theta=\{(i,j):i,j\in\{0,1\}\}$. Action
$R$ reports $i$. Action $L$ reports $j$ when $i=0$ and otherwise reports
an independent fair bit. Action $H$ reports $j$ when $i=1$ and otherwise
reports an independent fair bit. The candidate validation is
$\delta_{(i,j)}$ and the target is uniform on $\Theta$.

\begin{proposition}[Strict adaptive advantage]\label{prop:adaptive-gap}
With two training preparations, the unrestricted adaptive value is
$3/4$ and the best nonadaptive value is $5/8$. Nonadaptive policies here
include mixtures of schedules, with their random schedule retained.
No available single training action dominates all the others.
\end{proposition}
\begin{proof}
Take $R$ first, and then $L$ or $H$ according to the observed $i$.
The parameter is recovered exactly. The complement of its validation
atom has score $3/4$. No policy can exceed $\TV(Q,\delta_\theta)=3/4$
even with full information.

For the converse for fixed schedules, use the uniform prior. For a
pair of actions $a,b$ its Bayes value is explicitly
\begin{equation}\label{eq:adaptive-table-formula}
 \sum_{x,y,\theta}\left[
       \frac1{16}\sum_{\vartheta}P^a_\vartheta(x)P^b_\vartheta(y)
       -\frac14P^a_\theta(x)P^b_\theta(y)\right]_+.
\end{equation}
Substitution of the three displayed kernels gives
\[
\begin{array}{c|ccc}
 &R&L&H\\\hline
R&1/2&5/8&5/8\\
L&5/8&7/16&1/2\\
H&5/8&1/2&7/16
\end{array}.
\]
This is an exact finite sum, not a parameter-grid estimate. It bounds
any mixture of fixed schedules by $5/8$; an early stop can be padded
with an ignored observation. To attain the bound uniformly in $\theta$,
choose with equal probabilities the schedules $(R,L)$ and $(R,H)$.
When the chosen branch-specific action matches $i$, the parameter is
known and the score is $3/4$. Otherwise only $i$ is known, and the
complement of its two-atom fiber gives score $1/2$. Each fixed parameter
therefore has mean score $5/8$.

Action $R$ cannot distinguish parameters differing only in $j$.
Action $L$ cannot distinguish $(1,0)$ and $(1,1)$, and $H$ cannot
distinguish $(0,0)$ and $(0,1)$. An experiment which fails to distinguish
such a pair cannot garble to one that distinguishes it. Hence no action
is a common dominating row.
\end{proof}

If every binary report is additionally passed through an independent
binary symmetric channel of crossover $\epsilon$, the primitive
training rows change by at most $\epsilon$ and have full support for
$0<\epsilon<1/2$. Proposition~\ref{prop:controlled-stability}, applied
to both policy classes, leaves an adaptive advantage at least
$1/8-8\epsilon$, which is positive for $0<\epsilon<1/64$.
Thus the distinction is not confined to disjoint observation supports.
The finite-horizon belief recursion and its piecewise-linear structure
are classical \cite{SmallwoodSondik}. The contribution here is their
explicit integration with the frozen validation score, one fixed
adversarial parameter, outer certificates, and the resource conventions;
no priority for Bellman recursion itself is claimed.

\section{Exact minimax laws for a marked product family}
\label{sec:marked-exact}
The posterior-predictive dual can be solved explicitly for a continuous
family containing the collision target. Throughout this section, the
training and candidate validation law is
\[
 b(p,q)=\Ber(p)\otimes\Ber(q)\otimes\Ber(1/2),\qquad (p,q)\in[0,1]^2.
\]
In the order $000,001,110,111$, let $Q_\gamma$ have masses
\[
 ((1+\gamma)/4,(1-\gamma)/4,(1-\gamma)/4,(1+\gamma)/4),
\]
and zero mass elsewhere. Write $U_1(\gamma)=U_1(b,Q_\gamma)$, with
exactly the training--validation convention of \eqref{eq:selector}.
Set
\begin{equation}\label{eq:marked-constants}
 \gamma_*=1-1/\sqrt2,\quad s_\gamma=\sqrt{\gamma^2-2\gamma+5},\quad
 u_\gamma=\gamma-2+s_\gamma,\quad
 \tau_\gamma=\frac{2u_\gamma}{u_\gamma+2-\gamma}.
\end{equation}

\begin{theorem}[Exact marked minimax family]\label{thm:marked-exact}
For every $0\le\gamma\le\gamma_*$,
\begin{equation}\label{eq:marked-value}
 U_1(\gamma)=\frac{3+u_\gamma^2}{8}.
\end{equation}
An optimal response depends on the training word only through the report
count $K=X_1+X_2$. On the four target atoms its rows are
\begin{equation}\label{eq:marked-response}
 h_0=(\tau_\gamma,0,1,1),\qquad h_1=(1,1,1,1),\qquad
 h_2=(1,1,0,\tau_\gamma),
\end{equation}
and all responses are zero off these atoms. A least-favorable prior is
\begin{equation}\label{eq:marked-prior}
 \mu_\gamma=\tfrac12\delta_{(t_\gamma,t_\gamma)}
             +\tfrac12\delta_{(1-t_\gamma,1-t_\gamma)},\qquad
 t_\gamma=\frac{1-\sqrt{u_\gamma}}2.
\end{equation}
The assertion optimizes over all stochastic responses to the full marked
training word, not merely over count-based or symmetric responses.
\end{theorem}

\begin{proof}
First, $u_\gamma$ is the positive root of
\begin{equation}\label{eq:marked-root}
 u^2+(4-2\gamma)u-(1+2\gamma)=0.
\end{equation}
On the asserted interval, $0<u_\gamma<1$ and $0<\tau_\gamma<1$.
Indeed the polynomial is negative at zero and positive at one; also
$u_\gamma<2-\gamma$ because $\gamma<1/2$. Substitution of $1-2\gamma$
into the polynomial gives $8(\gamma^2-2\gamma+1/2)$, which is
nonnegative precisely on the present interval. Monotonicity on
$[0,\infty)$ therefore gives
\begin{equation}\label{eq:marked-range}
                    u_\gamma\le1-2\gamma.
\end{equation}

Put $a=(1-p)(1-q)$, $c=pq$, and $z=1-a-c$. The probabilities of the
three training counts are $a,z,c$. The independent training mark may
be discarded by the displayed construction. Let $G_\gamma(p,q)$ denote
its score. With $u=(p+q-1)^2$ and $v=(p-q)^2$, direct expansion gives
\begin{equation}\label{eq:marked-sos}
 G_\gamma(p,q)-\frac{3+u_\gamma^2}{8}
 =\frac{2-\tau_\gamma}{16}
       \big\{(u-v-u_\gamma)^2+4v\big\}.
\end{equation}
For clarity, this identity follows by substituting the three rows of
\eqref{eq:marked-response} into
\[
 a(Q_\gamma-b)h_0+z(Q_\gamma-b)h_1+c(Q_\gamma-b)h_2,
\]
using $a+c=(1+u-v)/2$, $a-c=-(p+q-1)$, and then eliminating
$u_\gamma^2$ by \eqref{eq:marked-root}. The coefficient of the resulting
square is $(2-\tau_\gamma)/16$ and that of $v$ is four times this
coefficient. Both are nonnegative. Equality holds at both points of
\eqref{eq:marked-prior}, since there $v=0$ and $u=u_\gamma$.
This proves the lower bound and identifies its equality set.

It remains to bound every response, including nonsymmetric ones. Use
\eqref{eq:marked-prior} in Theorem~\ref{thm:dual}. At $t=t_\gamma$ put
$a=(1-t)^2$, $c=t^2$, and $T=a+c=(1+u_\gamma)/2$. For the training
report pair $00$, the four dual coefficients, apart from a common
positive factor, are
\begin{equation}\label{eq:marked-coefficients}
 \left(Q_H T-\frac{a^2+c^2}{2},\quad
 Q_L T-\frac{a^2+c^2}{2},\quad Q_L T-ac,\quad Q_H T-ac\right),
\end{equation}
where $Q_H=(1+\gamma)/4$ and $Q_L=(1-\gamma)/4$.
The first coordinate is zero by \eqref{eq:marked-root}; the second
is $-\gamma T/2\le0$. The last two are positive: $T\ge1/2$,
$ac\le1/16$, and $\gamma<1/2$ imply $Q_LT>1/16\ge ac$.
Thus $(\tau_\gamma,0,1,1)$ maximizes this linear functional.
For the training pair $11$, reversing all report bits and the mark
reverses the four coordinates and gives the row $h_2$.

For either training pair $01$ or $10$, the two prior support points have
the same likelihood $t(1-t)$. The posterior is still equally weighted,
and the candidate predictive mass at each target atom is $T/4$.
Every one of the four coefficients is therefore nonnegative by
$Q_L-T/4=(1-2\gamma-u_\gamma)/8\ge0$.
The response $h_1$ is optimal. At off-diagonal validation atoms the target
mass vanishes and the coefficient is nonpositive. Finally, conditioning
on either training mark multiplies all these coefficients by $1/2$;
it creates no further optimizing choice. We have therefore optimized
every coordinate of the full marked-word response space.

The prior average of the displayed response equals its value at either
support point, namely $(3+u_\gamma^2)/8$, by \eqref{eq:marked-sos}.
Every other response has prior average at most this number and hence
has minimum over fixed $(p,q)$ at most this number. Together with the
lower bound this proves \eqref{eq:marked-value}, optimality of the
response, and least favorability of the prior.
\end{proof}

\begin{corollary}[The exact one-preparation collision value]
\label{cor:exact-one}
For $Q_*=Q_{1/4}$,
\begin{equation}\label{eq:exact-one}
 U_1(b,Q_*)=u_{\rm ex}:=\frac{85-7\sqrt{73}}{64},\quad
 u_{1/4}=\frac{\sqrt{73}-7}{4},\quad
 \tau_{1/4}=2-\frac{14}{\sqrt{73}}.
\end{equation}
If only the target changes to $Q$, then
$|U_1(b,Q)-u_{\rm ex}|\le\TV(Q,Q_*)$.
\end{corollary}
\begin{proof}
Substitute $\gamma=1/4$ in Theorem~\ref{thm:marked-exact}.
For a fixed trained response the target perturbation is at most total
variation, uniformly in $(p,q)$. Taking first the minimum and then the
maximum preserves this bound in both directions.
\end{proof}

\subsection{Complementary slackness and the event alphabet}
The tie in the first coordinate of \eqref{eq:marked-coefficients}
explains the randomization. Any value of that response coordinate is
Bayes-optimal for the two-point prior. The minimax choice makes the
uniform payoff stationary at its two support points; solving this
condition gives $\tau_\gamma$ and the square identity
\eqref{eq:marked-sos}. This derives the prior and response together,
rather than postulating that a numerically effective prior is optimal.

Operationally, for $K=0$ one draws once between the events with indicator
vectors $(0,0,1,1)$ and $(1,0,1,1)$, with probabilities $1-\tau_\gamma$
and $\tau_\gamma$. For $K=2$ one draws between $(1,1,0,0)$ and
$(1,1,0,1)$. For $K=1$ the event is $(1,1,1,1)$. The resulting event
identity is retained and used for both fresh validations. These are
exactly the five events in the two-preparation construction below.
The one-preparation theorem supplies their complementary-slackness
interpretation; it does not prove that the two-preparation selector is
optimal. The four-point rational prior retained below is a useful
rational certificate, not a least-favorable prior.

The exact response uses an atomic stochastic row. It does not receive
free persistent randomness. Replacing $\tau_\gamma$ by a dyadic number
within $2^{-b}$ changes each conditional response, and hence the score,
by at most $2^{-b}$. A sequential fair-bit implementation must retain its
bit prefix and event identity. None of these implementation costs is
needed by the deterministic two-preparation physical audit.

\begin{proposition}[The zero-preparation endpoint]\label{prop:marked-zero}
For $0\le\gamma\le1$, $U_0(b,Q_\gamma)=\gamma/2$.
\end{proposition}
\begin{proof}
The fixed event $\{000,111\}$ has score
$(1+\gamma)/2-(a+c)/2\ge\gamma/2$ because $a+c\le1$.
Conversely, the equally weighted prior on $b(0,0)$ and $b(1,1)$ has
predictive validation law uniform on the four target atoms. Its total
variation from $Q_\gamma$ is $\gamma/2$. Theorem~\ref{thm:dual} with
$N=0$ gives the matching upper bound.
\end{proof}

\section{A self-contained collision interface}
\label{sec:physical-bridge}
We give the preparation, mechanical estimate, observation law and
private simulator class needed for the physical specialization. The
following argument does not appeal to an unpublished physical lemma.
It retains the same positive-volume preparation and positive sensor
noise as the earlier collision construction.

\begin{definition}[The reset row--event architecture]
\label{def:physical-class}
A preparation supplies two sequential report slots and, afterwards, an
actual mark. The source reports in the two slots are constant. A private
width-one simulator has no retained information between those slots
and no persistent hidden seed. It may use fresh independent coins.
Its first emission is $X_1\sim\Ber(p)$; after receiving the current
gate $A\in\{0,1\}$, its second emission has a fixed Bernoulli kernel
with parameter $q_A$. The parameters $(p,q_0,q_1)$ are fixed throughout
the audit. The simulator does not observe the actual mark. All its
coins are independent of that mark. Different preparations reset all
unretained variables and are independent conditional on the parameters.

The auditor observes the first report, chooses a gate using its retained
state and that report, observes the second report, and then observes the
mark. During training it may choose rows and stop adaptively, using at
most the declared number of candidate preparations. It next freezes
one gate policy and one event, with any retained random choice included
in its register. It uses that same row--event continuation on one fresh
candidate preparation and one independent target preparation and returns
the difference of their event indicators, target minus candidate.
It may not reselect the event after either validation. Independent
target-only calibration is allowed but all retained calibration and
randomness are charged. Peak width counts every bit, mark, event and
validation-accumulator cut. Clocked schedules are declared externally;
an autonomous implementation must retain its phase. An unretained
transcript, seed or stopping-time value is not available to a continuation.
\end{definition}

In the all-on row the complete private law is $b(p,q_1)$, with an
independent fair actual mark. For lower bounds it suffices to restrict
the adversary to $q_0=0$, $q_1=q$. In this subfamily every deterministic
gate policy is the pushforward
\[
 T_g(x_1,x_2,w)=(x_1,g(x_1)x_2,w)
\]
of $b(p,q)$. The gate, when included in the transcript, is another
function of these same coordinates. Randomized gate policies are
covered by averaging their independent random tables. This proves
both the all-on product characterization for the upper construction
and the common-garbling subfamily used for the outer converse. The
width-one constraint is a simulator constraint; it is not an uncharged
restriction on the auditor's register.

\subsection{Preparation and elastic scattering}
Consider two labeled equal-mass hard spheres in $\mathbb R^3$, with no
walls and contact distance $d$. Write $r=q_1-q_2$, $g=v_1-v_2$,
$V=(v_1+v_2)/2$, and let the center of mass position range over any fixed
rational box of positive volume. Prepare $(B,W)\in\{0,1\}^2$ with
probability $5/16$ at $B=W$ and $3/16$ at $B\ne W$. Conditional on
$(B,W)$, use the uniform product preparation on the following boxes:
\begin{align}
 r&=(-a_0,(2B-1)b_0,c_0),&
 a_0&\in[29/10,31/10],& b_0&\in[9/20,11/20],\notag\\
 |c_0|&\le1/100,&g_x&\in[19/10,21/10],&
 |g_y|,|g_z|&\le1/100,\label{eq:local-packets}\\
 |V_x|,|V_y|&\le1/100,&&&|V_z-(2W-1)|&\le1/100.\notag
\end{align}
Thus the actual mark is $W=\mathbf1_{\{V_z>0\}}$, preserved by momentum
conservation. Set
\[
 h(z)=\operatorname{sgn}(z)\min\{1,(|z|-1/4)_+\}.
\]
This is a bounded $1$-Lipschitz detector, identically zero on the incoming
transverse-velocity range.

\begin{lemma}[Uniform collision signal]\label{lem:local-scattering}
For every preparation in \eqref{eq:local-packets} and every
$9/10\le d\le11/10$, exactly one collision occurs and its time lies in
$(4/5,3/2)$. Afterwards
\[
             (2B-1)v_{1,y}>1/2,\qquad (2B-1)h(v_{1,y})>1/4.
\]
For $0<d\le2/5$ no collision occurs on $[0,3]$ and the detector is zero
throughout that interval.
\end{lemma}
\begin{proof}
Before contact, $r(t)=r+tg$. At $t\le4/5$ its $x$ coordinate is at
most $-29/10+(4/5)(21/10)=-61/50$, so its norm exceeds $11/10$.
At $t=3/2$ its absolute coordinate bounds are $1/4$, $113/200$ and
$1/40$, whose squared sum is less than $(9/10)^2$. The first contact
therefore occurs in the claimed interval. At the free zero of the
$x$ coordinate, whose time is at most $31/19$, the transverse norm is
less than $9/10$. Hence the entry contact has $n_x<0$, where
$n=r(t_c)/d$. Its remaining coordinates give
\[
 |n_x|^2\ge1-\frac{(113/200)^2+(1/40)^2}{(9/10)^2}>(77/100)^2,
 \qquad (2B-1)n_y\ge87/220.
\]
It follows that $-g\cdot n>(19/10)(77/100)-2/100=1443/1000$.
The elastic rule gives
\[
 v_1^+=V+g/2-(g\cdot n)n,\qquad g^+=g-2(g\cdot n)n.
\]
Consequently
\[
 (2B-1)v_{1,y}^+>-3/200+(1443/1000)(87/220)>1/2.
\]
The definition of $h$ implies the signal bound. Since
$r(t_c)\cdot g^+>0$, squared separation increases strictly afterwards;
there is no second collision. Finally, for $d\le2/5$ and $0\le t\le3$,
$|(r+tg)_y|\ge9/20-3/100=21/50>d$. There is no contact, and the
incoming $|v_{1,y}|\le3/200$ makes $h$ zero.
\end{proof}

\subsection{Positive noise and the actual mark}
Let $\mathcal B_t$ be Brownian motion independent of the preparation and
set
\[
 Y_t=\int_0^t h(v_{1,y}(s))\,ds+\sigma\mathcal B_t,
                  \qquad 0<\sigma\le1/24.
\]
The two potential all-on reports are
\[
 Z_1=\mathbf1_{\{2(Y_2-Y_{3/2})>0\}},\qquad
 Z_2=\mathbf1_{\{2(Y_3-Y_{5/2})>0\}}.
\]
The actual second report is $AZ_2$. The gate changes acquisition, not
the mechanical path. The continuous path is implemented by the target
sensor; it is not made available to the private simulator or supplied
as an uncharged auditor buffer. Denote the law of $(Z_1,Z_2,W)$ by
$Q_{d,\sigma}$.

\begin{theorem}[Uniform physical presentation]\label{thm:local-physical}
For $9/10\le d\le11/10$ and $0<\sigma\le1/24$,
\begin{equation}\label{eq:local-tv}
                    \TV(Q_{d,\sigma},Q_*)<1/300.
\end{equation}
The same estimate holds after every declared feedback gate, including
its observable gate value. The all-on private product family and the
common-garbling subfamily in Definition~\ref{def:physical-class} apply
to this very experiment.
\end{theorem}
\begin{proof}
Every collision occurs before $3/2$. Conditional on the preparation,
the two normalized sensor windows are independent Gaussians with mean
$h(v_{1,y}^+)$ and standard deviation $\sqrt2\sigma\le1/12$.
Lemma~\ref{lem:local-scattering} makes their means have sign $2B-1$
and magnitude greater than $1/4$. Each potential report therefore
differs from $B$ with probability less than $\Phi(-3)$.
The Gaussian tail estimate
\[
 \Phi(-x)\le\frac{e^{-x^2/2}}{x\sqrt{2\pi}}
\]
gives $\Phi(-3)<1/600$: $e^{9/2}>80$ and $\sqrt{2\pi}>5/2$
suffice. Couple the ideal and noisy reports using the same preparation,
actual mark and gate random table. Except on a set of probability less
than $2\Phi(-3)<1/300$, both potential reports equal $B$ and the entire
gated transcript agrees. The ideal marked law is precisely $Q_*$,
by the prepared $(B,W)$ table. This proves \eqref{eq:local-tv} and the
feedback-uniform version. Finally, the source is blind, its actual mark
is unobserved by the simulator, and the gate is its only current input
at the second emission. The product and subfamily assertions follow
from the two independent private-emission kernels specified above.
\end{proof}

For $d\le2/5$, the two windows instead contain only independent centered
Gaussian noise. Their signs are independent fair bits, independent of
$W$, and a private width-one simulator reproduces every gated row
exactly. Thus the informative regime and its no-collision control use
the same detector, preparation convention and actual mark.


\section{The physical experiment and its sharp preparation boundary}
\label{sec:physical}
Take $\Omega=\{0,1\}^3$, with coordinates $(x_1,x_2,w)$. The all-on
private candidate is
\[
 b(p,q)=\Ber(p)\otimes\Ber(q)\otimes\Ber(1/2),
 \qquad 0\le p,q\le1.
\]
The mark $w$ is the actual mark, not an invented simulator bit. Put
$a=(1-p)(1-q)$ and $c=pq$. On the ordered atoms
\begin{equation}\label{eq:atoms}
       000,\quad001,\quad110,\quad111,
\end{equation}
let $\Qstar$ have masses $5/16,3/16,3/16,5/16$, and set it to zero
elsewhere. The candidate masses on these atoms are $a/2,a/2,c/2,c/2$.
A target $Q$ is allowed whenever $\TV(Q,\Qstar)\le\beta$.

The original feedback operation has the form
\begin{equation}\label{eq:gate}
 T_g(x_1,x_2,w)=(x_1,g(x_1)x_2,w),
 \qquad g:\{0,1\}\longrightarrow\{0,1\}.
\end{equation}
The selected gate is observable. It may depend on earlier preparations
and on the first current report. The target rows are pushforwards by
\eqref{eq:gate}. For the converse, restrict the adversary to the admissible
private subfamily whose second emission is $A X_2$, where $A$ is the
current gate and $X_2\sim\Ber(q)$ is independent of the first emission
and of the fair mark. Every row is then a causal garbling of $b(p,q)$.
For the construction only the all-on row is used, and every private
candidate in the original width-one class has the all-on product law
just specified. These two facts, proved for the microscopic interface in
Definition~\ref{def:physical-class} and Theorem~\ref{thm:local-physical}, explain
why using a subfamily for the lower bound does not narrow the upper
construction.

The admissible audit first trains, then selects one row--event pair, then
uses one fresh candidate and one fresh target validation. It reports
$D=\1_A(Y)-\1_A(X)$. The event and row cannot be reselected after
seeing a validation draw. All bit, mark, decision and validation cuts are
charged. Define $N_c^{\rm phys}(W,Q)$ to be the minimum number of candidate
training preparations among these original-interface auditors, with peak
register width at most $W$, whose score is uniformly strictly greater
than $c$. Randomization, feedback and stopping are allowed. If no such
audit exists, the value is $+\infty$.

\begin{theorem}[Two-preparation physical segment]\label{thm:physical}
Let $u_1,m_2$ be as in \eqref{eq:constants}. For every $\beta\ge0$,
$\TV(Q,\Qstar)\le\beta$, and
\[
       W\ge12,\qquad u_1+\beta\le c<m_2-\beta,
\]
one has $N_c^{\rm phys}(W,Q)=2$. The two-preparation audit is deterministic,
uses no acquired calibration and no random bits, and has the full profile
\begin{equation}\label{eq:profile}
 (1,2,3,3,4,5,5,10,12,10,10,7,3)
\end{equation}
from the initial cut through its twelve bit/mark updates. In particular,
its decision alphabet has five labels and its peak has twelve. The lower
bound remains valid with no memory cap.
\end{theorem}
The proof occupies the next three subsections. The interval in the
statement can be empty for large $\beta$; no assertion is then made.
For $\beta=1/300$, it contains $c=2/5$ strictly between its endpoints.

\subsection{An exact rational certificate against one preparation}
\begin{lemma}\label{lem:one}
Every original-interface auditor using at most one candidate training
preparation has worst-case score at most $u_1+\beta$, regardless of memory.
\end{lemma}
\begin{proof}
By Theorem~\ref{thm:outer}, it suffices to bound $U_1(b,\Qstar)$ and
then perturb the target. In \eqref{eq:dual} use the four fixed parameters
\[
 (p,q)=(1/8,1/8),(1/5,1/5),(4/5,4/5),(7/8,7/8)
\]
with weights $1/16,7/16,7/16,1/16$, respectively. This measure is not
asserted to be least favorable. Since $p=q$ on its support, the training
count $K=x_1+x_2$ is sufficient. The independent training mark and the
order of unequal bits contain no further information about this prior.

For $k=0,1,2$ and the four validation atoms in \eqref{eq:atoms}, the signed
coefficients in \eqref{eq:dual}, after grouping by $K$, are the entries of
\begin{equation}\label{eq:certificate}
 \frac1{40960000}
 \begin{pmatrix}
 -1977&-1775177&2170423&3943623\\
 1765554&191954&191954&1765554\\
 3943623&2170423&-1775177&-1977
 \end{pmatrix}.
\end{equation}
For example, the coefficient in column $j$ is
\[
 \sum_{i=1}^4\mu_i {2\choose k}p_i^k(1-p_i)^{2-k}
       \left(Q_{*,j}-\frac{(1-p_i)^2\1_{j\le2}
                                     +p_i^2\1_{j\ge3}}2\right).
\]
This identity proves every entry of \eqref{eq:certificate} by rational
arithmetic. All coefficients at off-diagonal validation atoms are
nonpositive. The sum of the positive entries in \eqref{eq:certificate}
is $4035777/10240000=u_1$. Each response coordinate is in $[0,1]$, so
this sum bounds its prior-averaged score and hence its minimum over the
four fixed candidates.

Replacing $\Qstar$ by $Q$ changes the score of any trained response by
at most $\beta$. The training law is a candidate law and is unchanged.
Consequently the same prior bounds every score by $u_1+\beta$. A procedure
with no training, or one that stops before taking its permitted sample,
is covered by ignoring the generated all-on observation.
\end{proof}

\subsection{A marked five-event selector}
Take two independent all-on training preparations. Let $S$ be the total
number of ones among their four report bits; discard their marks. Choose
the following events, represented by their indicators on
\eqref{eq:atoms} and with indicator zero at every off-diagonal atom:
\begin{equation}\label{eq:events}
\begin{array}{c|rrrr}
 S&000&001&110&111\\\hline
0&0&0&1&1\\
1&1&0&1&1\\
2&1&1&1&1\\
3&1&1&0&1\\
4&1&1&0&0
\end{array}
\end{equation}
Denote these indicators by $h_0,\ldots,h_4$ and report
$D=h_S(Y)-h_S(X)$ on independent validation draws.

\begin{lemma}[Exact uniform score of the selector]\label{lem:two}
Against the ideal target, the selector \eqref{eq:events} has uniform
minimum score $m_2=(12-3\sqrt3)/16$. Against $Q$, its score is at least
$m_2-\beta$.
\end{lemma}
\begin{proof}
The generating polynomial of $S$ is
\[
 ((1-p)+pz)^2((1-q)+qz)^2.
\]
Writing $g_s=(\Qstar-b(p,q))h_s$, its expected score is
$G(p,q)=\sum_{s=0}^4[z^s]((1-p)+pz)^2((1-q)+qz)^2 g_s$.
Substitution of the five rows in \eqref{eq:events} and collection of terms
with $x=p+q-1$, $y=p-q$, $u=x^2$ and $v=y^2$ gives
\begin{align}
 G(p,q)=H(u,v)
  ={}&-\frac{u^3}{32}+\frac{3u^2v}{32}+\frac{3u^2}{16}
       -\frac{3uv^2}{32}-\frac{uv}{4}-\frac{3u}{32}\notag\\
     &+\frac{v^3}{32}+\frac{v^2}{16}+\frac{15v}{32}+\frac7{16}.
                                                        \label{eq:poly}
\end{align}
This is a polynomial identity, not an interpolation over a grid. Since
$0\le u\le1$ and $v\ge0$ on the parameter square,
\[
 \partial_vH(u,v)=\frac{3(u-v)^2-8u+4v+15}{32}\ge\frac7{32}.
\]
Thus $H(u,v)\ge H(u,0)=(-u^3+6u^2-3u+14)/32$.
The derivative of the last expression is
$(-3u^2+12u-3)/32$. On $[0,1]$ it changes sign from negative to positive
at $u=2-\sqrt3$. Substitution gives $(12-3\sqrt3)/16$.
The minimizing values are feasible: take
$p=q=(1\pm\sqrt{2-\sqrt3})/2$. This proves exactness of the minimum
for this specified selector, not optimality of its score over all
selectors with two preparations. Finally, for each $s$, the target
expectation of $h_s\in[0,1]^\Omega$ changes by at most $\beta$.
Averaging over training proves the claimed robustness.
\end{proof}

\subsection{Causal compilation and the full register profile}
\begin{lemma}\label{lem:profile}
The selector and its two validations admit the deterministic profile
\eqref{eq:profile}. No selector, first-bit buffer or validation accumulator
is outside the charged register. An autonomous phase-tagged realization
uses at most $75$ states.
\end{lemma}
\begin{proof}
After each training report bit retain the count accumulated so far.
The two marks are read and discarded. The first six updates therefore
have widths $2,3,3,4,5,5$ after the initial singleton. At the decision
cut the five count values select the five events.

During candidate validation, after the first bit retain the event and
that bit, using at most ten states. After the second bit retain the event
and the function of the unread mark that gives its candidate-event
indicator. For events $0,1,2,3,4$, respectively, the possible functions are
\[
 \{0,1\},\quad\{0,1-w,1\},\quad\{0,1\},\quad
 \{0,w,1\},\quad\{0,1\}.
\]
There are twelve event--function pairs. After the candidate mark, retain
the event and its now known candidate indicator $c_0\in\{0,1\}$,
using at most ten states.

During target validation, after the first bit discard the event name
and retain $c_0$ and its residual function of the second bit and mark.
These residual functions belong to
\[
 \{0,\ x_2,\ (1-x_2)(1-w),\ 1-x_2,\ x_2w\}.
\]
This uses at most ten states. After the second bit, the residual
\emph{difference} as a function of the unread mark belongs to
\[
 \{-1,0,1,w,1-w,w-1,-w\},
\]
so seven states suffice. The mark update gives the three possible outputs
$-1,0,1$. Each transition depends only on its retained residual and the
current report. Thus it is a legal common observable row, not a family
of independently optimized hidden-history maps.

For completeness, equality of two residual functions implies equality of
their restrictions after either next bit, so identifying equal residuals
is closed under every subsequent update. The supplied transition tables
are obtained by these restrictions and can be reconstructed directly from
\eqref{eq:events}. Enumerating residuals is a convenience for reproducing
the tables, not the proof of their sufficiency. Finally, disjointly tag
the states at the thirteen cuts by their phase and use terminal
self-loops. The sum of the displayed widths is $75$. This implementation
charges the clock; it does not assert that $75$ is optimal.
\end{proof}

\begin{proof}[Proof of Theorem~\ref{thm:physical}]
Lemma~\ref{lem:one} excludes $N\le1$ for a strict score greater than
$c\ge u_1+\beta$, for every architecture and every memory budget.
Lemmas~\ref{lem:two} and \ref{lem:profile} supply a two-preparation
architecture of peak twelve with score at least $m_2-\beta>c$.
Larger width budgets can use the same machine.
\end{proof}

\begin{corollary}[The positive-noise collision family]\label{cor:collision}
For the original marked collision targets with
$9/10\le d\le11/10$ and $0<\sigma\le1/24$, the local microscopic
estimate $\TV(Q_{d,\sigma},\Qstar)<1/300$ implies
\[
             N_{2/5}^{\rm phys}(W,Q_{d,\sigma})=2
                         \qquad(W\ge12).
\]
There are four preparations in the complete audit: two candidate training,
one candidate validation, and one target validation.
\end{corollary}
\begin{proof}
Apply Theorem~\ref{thm:physical}. The lower-end margin is the exact rational
number
\[
 \frac25-u_1-\frac1{300}=\frac{78269}{30720000}>0.
\]
For the other endpoint, $\sqrt3<7/4$ implies
$m_2>27/64$ and $27/64-1/300>2/5$.
The physical total-variation estimate, including the actual mark and
positive sensor noise, is proved in Theorem~\ref{thm:local-physical}; the current deduction
does not replace those hypotheses by a zero-noise limit.
\end{proof}

\begin{remark}[The part of the frontier determined here]
The theorem evaluates the minimum preparation coordinate for $W\ge12$.
It does not determine $N_c^{\rm phys}(W,Q)$ for $3\le W\le11$, the
maximum score with two preparations, or the least peak width at $N=2$.
The earlier exact decision-width result $\kappa_{2/5}^{\rm rows}=3$
is unchanged: it allows longer training, whereas this shortest audit
uses five decision states. The older points $(399,40602)$,
$(320000,1604)$ and $(12800\cdot2^{400},12)$ remain valid constructions
with decision width three. They are dominated by $(2,12)$ in the
projection onto training preparations and peak width, but not necessarily
in a resource order that also constrains the decision alphabet separately.
\end{remark}

\begin{corollary}[The enlarged physical preparation segment]
\label{cor:improved-physical}
With the full architecture of Definition~\ref{def:physical-class},
\[
 N_c^{\rm phys}(W,Q)=2
 \quad\text{if}\quad W\ge12,\quad \TV(Q,Q_*)\le\beta,
 \quad u_{\rm ex}+\beta\le c<m_2-\beta.
\]
Here $u_{\rm ex}=(85-7\sqrt{73})/64$ is the exact one-preparation
all-on value, rather than the older rational upper bound.
\end{corollary}
\begin{proof}
Restrict the adversary to the common-garbling subfamily in
Definition~\ref{def:physical-class}. Theorem~\ref{thm:outer} and
Corollary~\ref{cor:exact-one} bound every at-most-one-preparation auditor,
even with unlimited memory, by $u_{\rm ex}+\beta$. The same deterministic
two-preparation machine in Lemmas~\ref{lem:two} and \ref{lem:profile}
has score at least $m_2-\beta$ against the full private candidate family.
These inequalities prove the assertion. The four-preparation total,
peak twelve and positive-noise specialization are unchanged.
\end{proof}

\section{An exact preparation--decision-memory law}
\label{sec:joint-law}
A second solvable family separates a genuine joint resource law from an
implementation bound. Here the memory coordinate is the number of states
at the training--validation cut. It is not identified with the bit-level
peak width of the collision audit.

Fix $m=2^d$, $0<\alpha<1$, and an unknown parameter
$\theta\in\{1,\ldots,m\}$. A fresh training preparation reports
$\theta$ with probability $\alpha$ and an erasure $\bot$ otherwise.
Conditional on $\theta$, the preparations are independent. Candidate
validation has law $\delta_\theta$ and target validation is uniform on
$\{1,\ldots,m\}$. At the decision cut the auditor retains at most
$K=2^k$ states, $0\le k\le d$. Its subsequent response may depend only
on that state, current validation observations and fresh independent
randomness. In particular, a calibration table, remembered seed or
stopping time outside the retained state is not available. Upstream and
downstream memory are otherwise unrestricted for the converse. Let
$V_N^{\rm er}(K)$ be the best worst-case expected signed validation score
with at most $N$ training preparations.

\begin{theorem}[Exact joint law]\label{thm:joint-law}
For every $N\ge0$ and $K=2^k\le m$,
\begin{equation}\label{eq:joint-law}
 V_N^{\rm er}(K)=\bigl(1-(1-\alpha)^N\bigr)\left(1-\frac1K\right).
\end{equation}
The converse covers stochastic encoders and arbitrary legal stopping.
The value is attained by a clocked machine with $K$ training states.
With atomic preparation and validation outcomes, the complete audit
has peak at most $\max\{2K,3\}$ and uses at most $N+2$ preparations.
\end{theorem}
\begin{proof}
Write $q=(1-\alpha)^N$ and pre-generate $N$ potential training reports.
A stopped auditor ignores the unused reports; this does not disclose
extra information to it. On the event that all $N$ reports are erasures,
its retained-state law $r=(r_j)_{j=1}^K$ is independent of $\theta$.
After padding any smaller register with unused states, its channel to
the decision register has the form
\begin{equation}\label{eq:erasure-channel}
    C(j\mid\theta)=q r_j+(1-q)S(j\mid\theta)
\end{equation}
for a stochastic channel $S$. This identity remains valid with randomized
stopping and randomized state updates.

Use the uniform prior on parameters as an upper certificate. The part
$q r_j$ has zero expected score for every response: averaging
$\delta_\theta$ under this prior gives the target itself. For the other
part, optimizing over all responses gives a convex function of $S$, a
sum of positive parts of linear functions. Every stochastic channel
is a convex combination of deterministic channels $j=g(\theta)$;
therefore its value is at most the largest deterministic-channel value.
This convexification is used only to bound a channel. It does not
presuppose that a mixture of width-$K$ machines is realizable at width
$K$.

If the fibers of $g$ have sizes $n_1,\ldots,n_K$, a nonempty fiber has
posterior uniform on that fiber, so its total variation from the uniform
target is $1-n_j/m$. The prior-averaged optimal score is
\[
 \sum_j\frac{n_j}{m}\left(1-\frac{n_j}{m}\right)
    =1-\sum_j(n_j/m)^2\le1-1/K.
\]
Empty fibers contribute zero. Cauchy--Schwarz proves the inequality.
This is an upper bound on the average score and hence on its minimum
over fixed parameters. Multiplying by $1-q$ proves the upper bound in
\eqref{eq:joint-law}.

For attainment, partition the parameter alphabet into $K$ equal groups
and let $g(\theta)$ denote its group. Initialize the register uniformly
on the $K$ group labels. On an erasure leave it unchanged; on a revealed
parameter overwrite it by $g(\theta)$. No extra ``revelation occurred''
flag is needed. At the decision cut its law is
$q\operatorname{Unif}[K]+(1-q)\delta_{g(\theta)}$.
For register $j$, freeze the event which is the complement of group $j$.
On a revelation its score is $1-1/K$. Under a uniformly random initial
label, the mean score is zero for each fixed $\theta$: one of the $K$
labels has score $1-1/K$ and the others have score $-1/K$.
This proves the exact worst-case value.

The $k$ initialization bits can be accumulated in the $K$-state register;
all earlier prefix counts are at most $K$. Thus no persistent seed lies
outside it. Candidate validation retains the group label and its event
indicator, using at most $2K$ states. Target validation produces one
of the three values $-1,0,1$. This proves the atomic-outcome peak bound.
The controller uses a declared external schedule; an autonomous stopping
implementation must additionally retain its phase. A phase-tagged
implementation is finite, but its cost is not included in $K$.
\end{proof}

\begin{corollary}[Thresholds and matching resource scales]
\label{cor:joint-scale}
For $K\ge2$ and $0\le c<1-1/K$, the least integer number of training
preparations giving a score strictly greater than $c$ is
\begin{equation}\label{eq:joint-threshold}
 N_c^{\rm er}(K)=
 \left\lfloor\frac{\log(1-c/(1-1/K))}{\log(1-\alpha)}\right\rfloor+1.
\end{equation}
If $c\ge1-1/K$, no finite $N$ suffices. For $0<\varepsilon<1$, achieving
$V_N^{\rm er}(K)\ge1-\varepsilon$ requires
\[
 K\ge\varepsilon^{-1},\qquad
 N\ge\frac{\log(1/\varepsilon)}{-\log(1-\alpha)}.
\]
Conversely, whenever an allowed $K\le m$ satisfies $K\ge2/\varepsilon$,
the choice
$N\ge\log(2/\varepsilon)/[-\log(1-\alpha)]$ suffices.
\end{corollary}
\begin{proof}
The strict threshold is equivalent to
$(1-\alpha)^N<1-c/(1-1/K)$. Since the denominator logarithm is
negative, the least integer satisfying this strict inequality is
\eqref{eq:joint-threshold}, including equality cases of the real
quotient. Also
\[
 1-V_N^{\rm er}(K)=q+(1-q)/K.
\]
This is at least $q$ and at least $1/K$, and is at most $q+1/K$.
These three inequalities prove the necessary and sufficient bounds.
\end{proof}

The exact law concerns an entire $(N,K,\alpha,m)$ family, not a single
finite diagnostic. Its mechanism is a parameter-independent erasure
component followed by the convex geometry of the decision channel.
It demonstrates how a preparation obstruction and a memory obstruction
can be matched in one theorem. It does not identify the least physical
peak width at two collision preparations. Encoding the $m$-symbol
reports bit by bit also introduces cuts not present in this atomic
model; the physical bit-level accounting is given separately.


\section{Quantitative certificates for a finite compatible architecture}
\label{sec:quantitative}
We return to an arbitrary finite architecture, now with a finite parameter
set. It can have hidden states, controlled experiments, more than two
hypotheses and registers of any size. Enumerate its distinct stochastic
rows by $r=1,\ldots,R$ and write $q_r\in\Delta_{s_r}$.
A reused autonomous row appears only once in this product. Assume each
execution uses at most $H$ controller rows, and uses row $r$ at most $d_r$
times. Delete unused rows and put $S=\max_rs_r$.
Fixed environmental kernels are not controller rows.

For an integer $m\ge1$, intersect $\Delta_{s_r}$ with the coordinate
boxes of side $1/m$. The resulting nonempty closed polytopes cover the
simplex, with overlaps permitted at their boundaries. Write $C$ for a
product of such row cells and choose a vertex $q^C_r$ in each factor.
A row-cell vertex has at most one coordinate not fixed at a box boundary;
consequently it has at most $v_r=s_r2^{s_r-1}$ vertices, all on the
$1/m$ grid. There are at most $m^{\sum_rs_r}$ product cells.

Write the vertices of row cell $C_r$ as $v_{r,j}$ and parametrize
$q_r=\sum_j\lambda_{r,j}v_{r,j}$ with $\lambda_r$ in a simplex.
A reward polynomial has degree at most $d_r$ in block $\lambda_r$.
Use the homogeneous simplex Bernstein basis of degree $d_r$ in that
block, obtaining
\begin{equation}\label{eq:bernstein}
 g_\theta(q(\lambda))=\sum_\alpha b_{\theta,C,\alpha}B_\alpha(\lambda),
 \qquad B_\alpha\ge0,\quad\sum_\alpha B_\alpha=1.
\end{equation}
Degree elevation includes the factors $\sum_j\lambda_{r,j}=1$.
No independence between different visits of a shared row is imposed on
its unknown row vector.

\begin{lemma}[Causal interpretation of the coefficients]\label{lem:urn}
Every coefficient in \eqref{eq:bernstein} satisfies
\[
       |b_{\theta,C,\alpha}-g_\theta(q^C)|\le HS/m.
\]
\end{lemma}
\begin{proof}
For each row $r$, independently prepare an ordered list of $d_r$ vertex
labels, sampled independently with law $\lambda_r$. At its $k$th visit,
use the stochastic row given by the $k$th vertex label in that list.
Before conditioning, the next label is independent of the previous
execution, so this experiment has exactly the row $q_r$ at every visit
and expected reward $g_\theta(q(\lambda))$. Unused labels may be sampled
after termination without altering the reward.

Condition on the counts of every vertex label in every list. The
probability of those counts is $B_\alpha(\lambda)$. Conditional on the
counts, each list is a uniformly ordered multiset. Its expected reward
is independent of $\lambda$ and is therefore the corresponding Bernstein
coefficient. This also proves \eqref{eq:bernstein} directly.

Each row used by this conditioned experiment is a vertex of $C_r$ or a
mixture of its remaining vertices. Every two vectors in $C_r$ have
coordinate differences at most $1/m$, hence total variation at most
$s_r/(2m)\le S/(2m)$. Couple the conditioned experiment to the one using
$q^C$, using the same environmental randomness while their executions
agree. At each actual controller update the probability of a first
row-level disagreement is at most $S/(2m)$. There are at most $H$ such
updates on the agreeing path. The terminal laws consequently differ in
total variation by at most $HS/(2m)$. A reward in $[-1,1]$ has expectation
difference at most twice that quantity. This proves the estimate. The
conditioned urn is used to analyze a coefficient; it is not declared to
be an admissible free-memory controller.
\end{proof}

\begin{theorem}[Explicit subdivision certificate]\label{thm:quantitative}
Define
\[
 \ell_C=\min_\theta g_\theta(q^C),\qquad
 u_C=\min_\theta\max_\alpha b_{\theta,C,\alpha},\qquad
 L_m=\max_C\ell_C,\qquad U_m=\max_Cu_C.
\]
For the exact all-stochastic value of the specified architecture,
\begin{equation}\label{eq:quantitative}
       L_m\le V(A)\le U_m\le L_m+HS/m.
\end{equation}
On each cell, a witnessing parameter $\theta_C$ supplies the nonnegative
coefficient certificate
\[
 u_C-g_{\theta_C}(q(\lambda))
       =\sum_\alpha(u_C-b_{\theta_C,C,\alpha})B_\alpha(\lambda).
\]
Its degree in row block $r$ is $d_r$, and its total degree is at most
$\sum_rd_r$. At most
$\prod_r{d_r+v_r-1\choose v_r-1}$ coefficients per parameter and cell
are needed. If the primitive data are rational, all coefficients and
bounds can be checked by rational arithmetic. Taking $m\ge HS/\varepsilon$
gives a certified interval of length at most $\varepsilon$.
\end{theorem}
\begin{proof}
The point $q^C$ is feasible, so $\ell_C$ is a lower bound on the value
restricted to $C$. For any $q\in C$, choose a convex representation of
its rows. Equation~\eqref{eq:bernstein} implies
\[
 \min_\theta g_\theta(q)\le
 \min_\theta\max_\alpha b_{\theta,C,\alpha}=u_C.
\]
Choose a parameter attaining the minimum defining $\ell_C$ and apply
Lemma~\ref{lem:urn}; this gives $u_C\le\ell_C+HS/m$.
Taking the maximum over the finite cover proves \eqref{eq:quantitative}.
Choose $\theta_C$ attaining the definition of $u_C$; all displayed
coefficients of its certificate are nonnegative. The degree and count
statements follow from the product Bernstein basis. Cell vertices are
rational grid points, and a finite execution sums finite products, proving
the rational statement.
\end{proof}

\begin{remark}
The number of cells and coefficients can be enormous. Neither polynomial
running time nor a uniform degree bound for the original sum-of-squares
hierarchy follows. Refinement changes the domain cover, not merely a
single global SOS degree. For a finite collection of architectures obeying
a specified budget, maxima of these intervals give outer intervals; an
unbounded architecture class needs a separate uniform argument such as
Theorem~\ref{thm:outer}. For compact parameter families, an effective net
and a verified modulus for the primitive kernels are additional inputs,
not consequences of compactness alone.
\end{remark}

The resource estimate also makes row precision explicit. With $m=2^b$,
the lower witnesses have dyadic rows and lose at most $HS2^{-b}$ in value.
Thus increasing row precision provides a quantified approximation for a
fixed architecture, while its random-bit implementation still has to be
charged. The physical theorem avoids this cost altogether because its
rows are deterministic.

\section{Autonomous noisy testing with a precision and time ledger}
\label{sec:autonomous}
This section supplies a substantive autonomous comparison rather than
merely identifying shared rows. The stationary optimum is classical
\cite{HC}; the proof and explicit finite-time implementation are recorded
to make the resource conventions unambiguous. They are not used as
novelty evidence for the physical theorem.

Let reports be independent, with law $\Ber(p)$ under hypothesis $-$ and
$\Ber(q)$ under hypothesis $+$, where $0<p<1/2$, $q=1-p$, and $r=q/p$.
A $W$-state autonomous machine has two time-independent stochastic
transition matrices $T_0,T_1$. Assume their induced chains
$P_-=qT_0+pT_1$ and $P_+=pT_0+qT_1$ are irreducible. A final decision
uses only the state, with optional fresh terminal randomization.

\begin{lemma}[Likelihood-range bound]\label{lem:range}
If equivalent probability laws $\nu_+,\nu_-$ on a finite set satisfy
\[
 \frac{\max_i\nu_+(i)/\nu_-(i)}
      {\min_i\nu_+(i)/\nu_-(i)}\le R,
\]
then $\TV(\nu_+,\nu_-)\le(\sqrt R-1)/(\sqrt R+1)$.
\end{lemma}
\begin{proof}
Write $L=d\nu_+/d\nu_-$, $a=\min L$ and $b=\max L$.
Since $\E_{\nu_-}L=1$, one has $a\le1\le b$. The chord bound for
$(L-1)_+$ gives
\[
 \TV(\nu_+,\nu_-)=\E_{\nu_-}(L-1)_+
 \le\frac{(1-a)(b-1)}{b-a}.
\]
For fixed $b/a$, maximizing the last expression gives $a=1/\sqrt{b/a}$;
substitution yields $(\sqrt{b/a}-1)/(\sqrt{b/a}+1)$. It is increasing
in $b/a$, proving the assertion. The identical-law case follows directly.
\end{proof}

\begin{theorem}[Autonomous stationary bound and an explicit approximation]
\label{thm:autonomous}
Put $D=W-1$, $W\ge2$. In the irreducible autonomous class the infimum
stationary minimax error is
\begin{equation}\label{eq:HC}
                         e_W=\frac1{1+r^D}.
\end{equation}
For every $\epsilon=2^{-b}$, $b\ge0$, there is a machine whose rows have
denominator dividing $2^{b+1}$, and whose minimax error at report time $n$
from any initial state is at most
\begin{equation}\label{eq:aut-time}
 e_W+\epsilon(W-2)
       +\exp\left\{-\left(\frac{p\epsilon}{2}\right)^D
                            \left\lfloor\frac nD\right\rfloor\right\}.
\end{equation}
The stationary lower bound is not a finite-horizon lower bound.
\end{theorem}
\begin{proof}
Every positive induced transition has its likelihood ratio in
$[r^{-1},r]$, and the supports of the two induced matrices are identical.
The Markov chain tree formula writes each stationary weight as a sum of
products over directed spanning trees rooted at that state. Each product
has $D$ transitions; its ratio lies in $[r^{-D},r^D]$. The ratio of the
two unnormalized stationary weights lies in the same interval. Normalizing
multiplies all state likelihood ratios by a common number, so their
maximum-to-minimum ratio is at most $r^{2D}$. Lemma~\ref{lem:range}
therefore bounds stationary total variation by $(r^D-1)/(r^D+1)$.
The equal-prior Bayes error is at least $1/(1+r^D)$, and a minimax error
cannot be smaller than that Bayes error. This is the equal-prior binary
specialization of the Hellman--Cover bound.

For the upper approximation, take states $i=0,\ldots,D$, weights
$c_0=c_D=1$, and $c_i=\epsilon$ at interior states. On report $1$, move
from $i$ to $i+1$ with probability
$a_i=\frac12\min(1,c_{i+1}/c_i)$, if that neighbor exists; otherwise
stay. On report $0$, move to $i-1$ with probability
$b_i=\frac12\min(1,c_{i-1}/c_i)$, if it exists; otherwise stay.
All other probability is assigned to staying. Detailed balance gives
\[
 \pi_+(i)=\frac{c_ir^i}{\sum_jc_jr^j},\qquad
 \pi_-(i)=\frac{c_ir^{-i}}{\sum_jc_jr^{-j}}.
\]
Reflection $i\mapsto D-i$ exchanges these laws. Decide $+$ for $i>D/2$,
$-$ for $i<D/2$, and use a fair terminal decision at $i=D/2$ when this
state exists. The two errors agree. Conditional on being at an endpoint,
the error is exactly $e_W$. The total interior stationary mass is at most
$\epsilon(W-2)$, since each $r^i\le r^D$ and the denominator is at
least $r^D$ under $+$, with the reflected argument under $-$.
This proves the first two terms of \eqref{eq:aut-time}.

Every existing neighbor transition under either hypothesis has probability
at least $p\epsilon/2$, and every self-loop has probability at least
$1/2$. From any state, move down to zero and then stay until $D$ steps
have elapsed. This yields the common minorization
$P_\pm^D(i,\{0\})\ge\delta=(p\epsilon/2)^D$.
Subtracting that common atom gives contraction in total variation by
$1-\delta$ every $D$ steps. Comparison with the stationary law yields
$\TV(P_\pm^n(i,\cdot),\pi_\pm)\le
(1-\delta)^{\lfloor n/D\rfloor}\le e^{-\delta\lfloor n/D\rfloor}$.
A bounded decision changes error by at most this amount, proving
\eqref{eq:aut-time}. Finally, letting $n$ tend to infinity and then
$\epsilon$ to zero proves that the stationary infimum equals
\eqref{eq:HC}. All transition probabilities belong to
$\{0,1,1/2,\epsilon/2,1-1/2,1-\epsilon/2\}$.
\end{proof}

The $W$ states in this theorem refer to an atomic stochastic update after
each report. A sequential fair-bit implementation must also retain the
report, the current state, a bit counter and a flag. Each nontrivial move
probability is $1/2$ or $2^{-(b+1)}$; test for a string of zero bits of
the required length, and pad to $b+1$ bits if a uniform report cycle is
required. At most $4W(b+2)$ states suffice for this implementation,
including its bit-phase counter and a ready phase. This conservative
bound is an implementation claim, not an optimal precision--state
tradeoff. The time index $n$ in \eqref{eq:aut-time} is an observation time
at which an external analyst reads the autonomous machine, not a free
input clock used by its transitions. The adversarial model of
\cite{MP} is not covered by the independent-report proof above.

\section{Physical consumers and the first-principles connection}
\label{sec:consumers}
The physical preparation theorem needs no estimate of the target law as
a stored input: the five events are fixed and their uniform validity uses
a proved neighborhood of $\Qstar$. This removes calibration memory from
this particular consumer without asserting that target acquisition is
unnecessary in an arbitrary comparison problem. The general acquired
calibration theorem, with its full counter register, remains in the
complete companion.

\begin{corollary}[A raw-interface confidence construction]\label{cor:confidence}
Let $m=m_2-\beta>0$ and $0<c<m$. Repeat the full four-preparation audit
independently $R$ times and write $\overline D_R$ for the mean score.
Under every private candidate,
\[
 \Pp(\overline D_R\le c)\le\exp\{-R(m-c)^2/2\}.
\]
Under the null in which the independent candidate and target validation
laws agree, with the same event chosen before both validations,
\[
 \Pp(\overline D_R>c)\le\exp\{-Rc^2/2\}.
\]
A clocked implementation uses $4R$ preparations and at most
$12(2R+1)$ persistent states, with no calibration and no random bits.
\end{corollary}
\begin{proof}
The independent block scores lie in $[-1,1]$ and have mean at least $m$
under a private candidate. Hoeffding's inequality gives the first bound.
Under the null, conditional on the selected event, the two independent
validation indicators have the same mean. Each block therefore has mean
zero, and the same inequality gives the second bound. Retain the integer
sum in $\{-R,\ldots,R\}$ along with the current twelve-state block.
At a block boundary add its terminal score and reset its register.
The deterministic external schedule is part of the clocked convention.
\end{proof}

This is an upper construction on the raw marked interface. Its proof
does not transfer a lower bound from a deliberately reduced Bernoulli
interface, and no optimal confidence-memory frontier is claimed.

\begin{proposition}[Target transport without a training multiplier]
\label{prop:transport}
For the two-preparation audit, changing only the target from $Q$ to $Q'$
changes its expected score by at most $\TV(Q,Q')$. Its candidate training
law, register profile and event selector remain unchanged.
\end{proposition}
\begin{proof}
Condition on the candidate training word. The expectation difference is
$(Q-Q')h_S$, whose absolute value is at most $\TV(Q,Q')$. Average over
that unchanged training law.
\end{proof}
In particular, every common-preparation transport estimate verified for
the original positive-noise collision family can be inserted here
without multiplying it by a long training horizon. If candidate kernels
also change, their training and validation defects must be added by a
causal coupling; the target-only statement does not remove those terms.

The first-principles route is now concrete. The positive experimental
kernels and actual reset preparations give the likelihoods in
\eqref{eq:dual}; these generate posterior continuation states. The
validation-test quotient gives the positive parts in the statistical
dual. A prior then supplies an architecture-uniform impossibility
certificate. On the constructive side, the five-event decision is
compiled by taking its future residual functions, which are closed under
observation shifts. These are compatible causal states and determine the
charged profile. Thus the preparation law, predictive continuation and
finite-resource realization are used in the same upper/lower theorem.
This is not a deduction of geometry or nonlinear operator compactness
from the existence of finite stochastic rows.

\section{Dependencies, retained results and remaining distinctions}
The exact three-report binary noisy frontier, arbitrary-horizon
binary-register minimax recursion and compatible Bayesian interval
recursion remain as proved in the complete companion. Its exact
odd-horizon majority realization is also retained; it is not a new result
of this revision. Section~\ref{sec:quantitative} applies to registers
of any finite width, controlled rows and any finite parameter set, but
provides certified intervals rather than a closed minimax recursion for
all such classes. The physical theorem supplies a matched preparation
segment in a continuous two-parameter family with overlapping stochastic
reports; its converse is not a finite parameter-grid check.

The stationary autonomous result in Section~\ref{sec:autonomous} has an
explicit formula, but its stationary lower bound must not be applied to
a finite-horizon clocked machine. Its finite-time upper bound is purposely
conservative and exposes the slowing caused by finite transition
precision. An arbitrary compact-family presentation still requires
additional effective information to construct a kernel net. A subdivided
positive-basis degree bound must not be called a degree bound for a
global SOS hierarchy.

The physical result is a genuine consumer of the experimental and
continuation root. It does not make the independent A2 geometric proofs
dependent on this root. The original B4 normalized-resolvent, range,
compactness and control-transfer obligations, and the broad C2
strict-dual, form, rigidity and optional-projection obligations, are not
proved by our finite-state theorem. Their statements, derivations and
current status remain in the retained historical pipeline. No aggregate
closure is inferred from preserving their source files.

Finally, the rational certificate \eqref{eq:certificate}, polynomial
identity \eqref{eq:poly}, and residual transitions are independently
reproducible. The program in the accompanying source checks their exact
identities and includes deliberately invalid controls. These checks are
not substitutes for the proofs above. Neither compilation, preservation
of earlier pages nor a successful finite test establishes novelty or an
editorial decision.

\begin{thebibliography}{99}
\bibitem{SmallwoodSondik} R. D. Smallwood and E. J. Sondik,
The optimal control of partially observable Markov processes over a finite horizon,
\emph{Operations Research} \textbf{21} (1973), 1071--1088.
\bibitem{Blackwell} D. Blackwell, Equivalent comparisons of experiments,
\emph{Ann. Math. Statist.} \textbf{24} (1953), 265--272.
\bibitem{Sion} M. Sion, On general minimax theorems,
\emph{Pacific J. Math.} \textbf{8} (1958), 171--176.
\bibitem{MM} J. M\"uller and G. Mont\'ufar, The geometry of memoryless
stochastic policy optimization in infinite-horizon POMDPs,
arXiv:2110.07409 (2021; revised versions).
\bibitem{Putinar} M. Putinar, Positive polynomials on compact semi-algebraic
sets, \emph{Indiana Univ. Math. J.} \textbf{42} (1993), 969--984.
\bibitem{HC} M. E. Hellman and T. M. Cover, Learning with finite memory,
\emph{Ann. Math. Statist.} \textbf{41} (1970), 765--782.
\bibitem{HCrandom} M. E. Hellman and T. M. Cover, On memory saved by
randomization, \emph{Ann. Math. Statist.} \textbf{42} (1971), 1075--1078.
\bibitem{FH} R. A. Flower and M. E. Hellman, Hypothesis testing with
finite memory in finite time, \emph{IEEE Trans. Inform. Theory}
\textbf{18} (1972), 429--431.
\bibitem{CFH} T. M. Cover, M. A. Freedman and M. E. Hellman, Optimal
finite memory learning algorithms for the finite sample problem,
\emph{Information and Control} \textbf{30} (1976), 49--85.
\bibitem{MP} M. A. Managoli and V. M. Prabhakaran, Memory constrained
adversarial hypothesis testing, arXiv:2605.12063v1 (May 12, 2026).
\bibitem{Norberg} E. Norberg, Comparison of statistical experiments with
filtered probability spaces, \emph{Statistics \& Risk Modeling}
\textbf{20} (2002), 1--28, doi:10.1524/strm.2002.20.14.1.
\bibitem{Companion} Q. Qi, General Theta Foundations I: Compatible
statistical experiments and resource frontiers, revision 23,
September 24, 2026, complete mathematical companion, source commit
\texttt{8a6043aca644043bd805668185eb227468f21170}; packaged commit
\texttt{e926b783b643b3e7f275a18105528927b76cc97b}.
The 96-page article and 649-page development are reproduced without
alteration after the present article in the two complete volumes.
\end{thebibliography}
\end{document}
